\section{Bab 1}
\subsection{Perumusan Dasar Negara}
\begin{itemize}
\item \textbf{1 Maret 1945: }Pembentukan BPUPKI.
\item \textbf{1 Juni 1945: }Sidang pertama pembentukan panita kecil (Memeriksa dan melaportkan sidang pleno).
\end{itemize}

\subsubsection{Anggota Panitia Sembilan}
\begin{enumerate}
\item IR. Soekarno (Ketua)
\item Ki bagoes Hadikoesoemo
\item K.H Wachid Hasyim
\item Moh Yamin
\item M. Soetardjo Kartuhadikoesomo
\item Mr. A.A Maramis
\item R. Otto iskandardinata
\item Moh hatta
\end{enumerate}

Panitia sembilan (Mempertemukan pandangan 2 golongan yang berbeda $\rightarrow$ Kebangsaan - Islam).

\begin{itemize}
\item \textbf{1 Juni 1945: }Kesepakatan preambul atau rumusan pancasila.
\item \textbf{22 Juni 1945: }Panitia sembilan $\rightarrow$ Menyetujui rancangan UUD 1945.
\item \textbf{10-17 Juli 1945: }Pelaporan membahas tentang piagam jakarta
\item \textbf{11 Juli 1945: }Tanggapan dari J. Latuharhary ``Ia tidak setuju atas tujuh kata: Ketuhanan dengan mewajibkan menjalankan syariat islam bagi peneluk pemeluknya.''
\end{itemize}

Akhir sidang BPUPKI $\rightarrow$ Berhasil menyusun $\rightarrow$ Dasar negara \textit{(groundnorm)} yang menjiwai atura dasar \textit{(groundgesetze)}

\begin{itemize}
\item \textbf{7 Agustus 1945: }BPUPKI dibubarkan. \\
  Pengesahan dasar negara $\rightarrow$ Menyerahnya jepang $\rightarrow$ Jatuhnya bom atom di kota hiroshima (6 Agustus 1945).
\item \textbf{7 Agustus 1945: }Pembentukan ``Dokuritsu junbi linkai'' $\rightarrow$ PPKI.
\item \textbf{9 Agustus 1945: }Sepulang dari saigon $\rightarrow$ IR. Soekarno, Moh Hatta, KRT Rdjiman Wedyodiningrat $\rightarrow$ Membentuk PPKI (21 orang).
\item \textbf{14 Agustus 1945: }Jepang menyerah kepada sekutu $\rightarrow$ Ada kekosongan kekuasaan.
\item \textbf{15 Agustus 1945: }IR. Soekarno, Moh Hatta, KRT Radjiman $\rightarrow$ Didesak golongan muda.
\end{itemize}

Pembacaan proklamasi awalnya $\rightarrow$ Di lapangan IKADA berubah menjadi di JL. Pegangsaan Timur No. 56 (JL Proklamasi 1) pada jum'at 17 Agustus 1945 10:00 WIB.


Rancangan UUD $\rightarrow$ ``Piagam Jakarta''.
